\documentclass{beamer}

% Theme choice
\usetheme{Madrid} % You can change to e.g., Warsaw, Berlin, CambridgeUS, etc.

% Encoding and font
\usepackage[utf8]{inputenc}
\usepackage[T1]{fontenc}

% Graphics and tables
\usepackage{graphicx}
\usepackage{booktabs}

% Code listings
\usepackage{listings}
\lstset{
basicstyle=\ttfamily\small,
keywordstyle=\color{blue},
commentstyle=\color{gray},
stringstyle=\color{red},
breaklines=true,
frame=single
}

% Math packages
\usepackage{amsmath}
\usepackage{amssymb}

% Colors
\usepackage{xcolor}

% TikZ and PGFPlots
\usepackage{tikz}
\usepackage{pgfplots}
\pgfplotsset{compat=1.18}
\usetikzlibrary{positioning}

% Hyperlinks
\usepackage{hyperref}

% Title information
\title{Introduction to Libraries: NumPy}
\author{Your Name}
\institute{Your Institution}
\date{\today}

\begin{document}

\frame{\titlepage}

\begin{frame}[fragile]
    \frametitle{Introduction to NumPy}
    \begin{block}{What is NumPy?}
        NumPy, short for "Numerical Python," is a powerful open-source library in Python designed for numerical computing. 
        It provides support for large, multi-dimensional arrays and matrices, along with a collection of mathematical functions to operate on these structures efficiently.
    \end{block}
\end{frame}

\begin{frame}[fragile]
    \frametitle{Significance of NumPy in Numeric Data Analysis}
    \begin{enumerate}
        \item \textbf{Performance}
        \begin{itemize}
            \item Built on C and Fortran, NumPy offers faster operations on arrays than standard Python lists.
            \item \textit{Vectorization}: Perform operations on entire arrays without loops, improving performance and reducing complexity.
        \end{itemize}
        
        \item \textbf{Ease of Use}
        \begin{itemize}
            \item Intuitive and concise syntax for straightforward mathematical operations.
            \item Built-in functions simplify complex operations (e.g., linear algebra, statistical analysis).
        \end{itemize}
        
        \item \textbf{Multidimensional Arrays}
        \begin{itemize}
            \item Unlike Python lists, NumPy’s \texttt{ndarray} can handle vast amounts of data efficiently in various dimensions.
        \end{itemize}
        
        \item \textbf{Mathematical Functions}
        \begin{itemize}
            \item Rich library enabling complex calculations, transformations, and statistical analyses effortlessly.
        \end{itemize}
    \end{enumerate}
\end{frame}

\begin{frame}[fragile]
    \frametitle{Key Features of NumPy}
    \begin{itemize}
        \item \textbf{ndarray}: A fast and flexible container for large datasets.
        \begin{lstlisting}[language=Python]
import numpy as np
# Creating a 1D array
array_1d = np.array([1, 2, 3, 4, 5])
        \end{lstlisting}

        \item \textbf{Broadcasting}: Mechanism allowing operations on arrays of different shapes.

        \item \textbf{Mathematical Operations}:
        \begin{itemize}
            \item Basic: Addition, subtraction, multiplication, etc.
            \item Advanced: Dot product, matrix operations, Fourier transforms, etc.
        \end{itemize}

        \item \textbf{Linear Algebra}: Functions for matrix multiplication and eigenvalue computations.
    \end{itemize}
\end{frame}

\begin{frame}[fragile]
    \frametitle{Example of NumPy Usage}
    \begin{lstlisting}[language=Python]
import numpy as np

# Creating a 2D array (matrix)
matrix = np.array([[1, 2, 3], 
                   [4, 5, 6], 
                   [7, 8, 9]])

# Perform operation: Transpose of the matrix
transpose = matrix.T

print("Original Matrix:\n", matrix)
print("Transposed Matrix:\n", transpose)
    \end{lstlisting}
    \begin{block}{Explanation}
        In this example, we create a 2D array (matrix) and compute its transpose using NumPy's \texttt{T} attribute. 
        This showcases the ease and efficiency of handling matrix operations.
    \end{block}
\end{frame}

\begin{frame}[fragile]
    \frametitle{Summary}
    \begin{block}{}
        NumPy is an essential library for numeric data analysis in Python. Its speed, flexibility, and rich functionality make it foundational for data scientists, engineers, and researchers.
        \begin{itemize}
            \item Mastering NumPy lays the groundwork for more complex data manipulation and analysis tasks.
        \end{itemize}
    \end{block}
    \begin{block}{Tip}
        Remember to practice using NumPy to become comfortable with its functionalities, as it will significantly enhance your data analysis capabilities in Python!
    \end{block}
\end{frame}

\begin{frame}[fragile]
    \frametitle{Installation and Setup - Pre-Requisites}
    \begin{itemize}
        \item \textbf{Python:} Ensure that you have Python installed on your system. Python 3.6 or later is recommended.
        \item \textbf{Package Manager:} Familiarize yourself with \texttt{pip}, the package manager for Python, which will be used to install NumPy.
    \end{itemize}
\end{frame}

\begin{frame}[fragile]
    \frametitle{Installation and Setup - Development Environment}
    Choose an Integrated Development Environment (IDE) that supports Python. Here are a few popular options:
    \begin{itemize}
        \item \textbf{Jupyter Notebook:} Ideal for interactive coding and data visualization.
        \item \textbf{PyCharm:} A powerful IDE with extensive features for Python development.
        \item \textbf{VS Code:} Lightweight and customizable, suitable for general coding and Python.
    \end{itemize}
\end{frame}

\begin{frame}[fragile]
    \frametitle{Installation and Setup - Installing NumPy}
    \textbf{Installing NumPy:}
    \begin{enumerate}
        \item Open your terminal (Command Prompt, PowerShell, or terminal emulator).
        \item Type the following command and hit Enter:
        \begin{lstlisting}
pip install numpy
        \end{lstlisting}
    \end{enumerate}
    
    \textbf{For Linux:}
    \begin{enumerate}
        \item Open the terminal.
        \item Run the same command:
        \begin{lstlisting}
pip install numpy
        \end{lstlisting}
    \end{enumerate}
\end{frame}

\begin{frame}[fragile]
    \frametitle{Installation and Setup - Verifying Installation}
    \textbf{Verifying the Installation:}
    After installation, verify if NumPy is successfully installed by running Python in your terminal:
    \begin{lstlisting}
import numpy as np
print(np.__version__)
    \end{lstlisting}
    If NumPy is installed correctly, this will print the version number of NumPy.
\end{frame}

\begin{frame}[fragile]
    \frametitle{Installation and Setup - Common Issues}
    \textbf{Common Installation Issues:}
    \begin{itemize}
        \item \textbf{Permission Errors:} Use \texttt{pip install --user numpy} to install it for your user account.
        \item \textbf{Outdated Pip:} If you experience issues, ensure that pip is up to date by running:
        \begin{lstlisting}
pip install --upgrade pip
        \end{lstlisting}
    \end{itemize}
\end{frame}

\begin{frame}[fragile]
    \frametitle{Installation and Setup - Key Points}
    \begin{itemize}
        \item Ensure Python 3.6+ is installed.
        \item Familiarize yourself with your chosen IDE, especially how to run Python scripts or notebooks.
        \item Installation success can be verified by checking the NumPy version.
        \item Consider using virtual environments for package management to avoid conflicts.
    \end{itemize}
\end{frame}

\begin{frame}[fragile]
    \frametitle{Installation and Setup - Example Code Snippet}
    \textbf{Example Code Snippet:} Confirm functionality post-installation.
    \begin{lstlisting}
import numpy as np

# Create a NumPy array
arr = np.array([1, 2, 3, 4, 5])
print("Created NumPy Array:", arr)
    \end{lstlisting}
\end{frame}

\begin{frame}[fragile]
    \frametitle{Key Features of NumPy - Overview}
    \begin{block}{Summary}
        This presentation will discuss three essential features of NumPy:
        \begin{itemize}
            \item ndarrays (N-dimensional arrays)
            \item Broadcasting
            \item Vectorization
        \end{itemize}
        Understanding these concepts is foundational for performing advanced data manipulations and calculations in scientific computing and data analysis.
    \end{block}
\end{frame}

\begin{frame}[fragile]
    \frametitle{Key Features of NumPy - ndarrays}
    \begin{block}{1. ndarrays (N-dimensional array)}
        \textbf{Definition:} The core feature of NumPy is the \texttt{ndarray}, a powerful, flexible container for large data sets. It efficiently handles multi-dimensional data and allows for operations on this data.

        \textbf{Example:}
        \begin{lstlisting}[language=Python]
import numpy as np

# Creating a 1D array
array_1d = np.array([1, 2, 3])

# Creating a 2D array
array_2d = np.array([[1, 2, 3], [4, 5, 6]])

# Shape of the array
print(array_2d.shape)  # Output: (2, 3)
        \end{lstlisting}

        \textbf{Key Points:}
        \begin{itemize}
            \item Supports various data types (int, float, etc.).
            \item Faster processing than Python lists.
            \item Convenient for data manipulation.
        \end{itemize}
    \end{block}
\end{frame}

\begin{frame}[fragile]
    \frametitle{Key Features of NumPy - Broadcasting}
    \begin{block}{2. Broadcasting}
        \textbf{Definition:} Broadcasting allows NumPy to perform element-wise operations on arrays of different shapes by automatically expanding the smaller array to match the larger array.

        \textbf{Example:}
        \begin{lstlisting}[language=Python]
a = np.array([1, 2, 3])      # Shape (3,)
b = np.array([[10], [20]])   # Shape (2, 1)

result = a + b
print(result)
# Output:
# [[11 12 13]
#  [21 22 23]]
        \end{lstlisting}

        \textbf{Key Points:}
        \begin{itemize}
            \item Reduces the need for explicit loops.
            \item Enhances performance through vectorized operations.
            \item Offers flexibility in array shapes for calculations.
        \end{itemize}
    \end{block}
\end{frame}

\begin{frame}[fragile]
    \frametitle{Key Features of NumPy - Vectorization}
    \begin{block}{3. Vectorization}
        \textbf{Definition:} Vectorization converts operations that require loops into operations that apply to entire arrays, greatly speeding up computations.

        \textbf{Example:}
        \begin{lstlisting}[language=Python]
# Using broadcasting and vectorization to compute squares
x = np.array([1, 2, 3, 4])

# Squaring the array without a loop
squares = x ** 2
print(squares)  # Output: [ 1  4  9 16]
        \end{lstlisting}

        \textbf{Key Points:}
        \begin{itemize}
            \item Provides clean, concise code.
            \item Leverages NumPy's optimized and compiled C code for speed.
            \item Ideal for mathematical and data analysis applications.
        \end{itemize}
    \end{block}
\end{frame}

\begin{frame}[fragile]
    \frametitle{Creating Arrays - Overview}
    In this slide, we will explore how to create arrays using the NumPy library, a powerful tool for numerical computations in Python. 
    \begin{itemize}
        \item Different methods of creating arrays
        \item Manipulating their structures for desired operations
    \end{itemize}
\end{frame}

\begin{frame}[fragile]
    \frametitle{Creating Arrays - Why Use NumPy Arrays?}
    \begin{itemize}
        \item \textbf{Efficiency}: NumPy arrays are more memory-efficient than Python lists.
        \item \textbf{Performance}: Operations on NumPy arrays are faster due to optimized C and Fortran underpinnings.
        \item \textbf{Multidimensional}: Arrays can have multiple dimensions, essential for mathematical computations in data science and machine learning.
    \end{itemize}
\end{frame}

\begin{frame}[fragile]
    \frametitle{Creating NumPy Arrays - Import and Examples}
    \begin{block}{1. Importing NumPy}
        Before we can start creating arrays, we need to import the NumPy library:
        \begin{lstlisting}
import numpy as np
        \end{lstlisting}
    \end{block}
    
    \begin{block}{2. Creating Arrays from Lists}
        You can create a NumPy array from a Python list using \texttt{np.array()}:
        \begin{lstlisting}
# Creating a 1D array
one_d_array = np.array([1, 2, 3, 4, 5])
# Creating a 2D array (matrix)
two_d_array = np.array([[1, 2, 3], [4, 5, 6]])
        \end{lstlisting}
        Key Point: Dimensionality is inferred from nested lists.
    \end{block}
\end{frame}

\begin{frame}[fragile]
    \frametitle{Creating NumPy Arrays - Functions and Reshaping}
    \begin{block}{3. Creating Arrays with Functions}
        \begin{itemize}
            \item \texttt{np.zeros(shape)}: Creates an array filled with zeros.
            \begin{lstlisting}
zeros_array = np.zeros((3, 4))  # 3 rows and 4 columns
            \end{lstlisting}

            \item \texttt{np.ones(shape)}: Creates an array filled with ones.
            \begin{lstlisting}
ones_array = np.ones((2, 3))  # 2 rows and 3 columns
            \end{lstlisting}

            \item \texttt{np.arange(start, stop, step)}: Creates an array with a range of values.
            \begin{lstlisting}
range_array = np.arange(0, 10, 2)  # [0, 2, 4, 6, 8]
            \end{lstlisting}
            
            \item \texttt{np.linspace(start, stop, num)}: Creates an array of evenly spaced values.
            \begin{lstlisting}
linspace_array = np.linspace(0, 1, 5)  # [0. , 0.25, 0.5, 0.75, 1.]
            \end{lstlisting}
        \end{itemize}
    \end{block}
    
    \begin{block}{4. Reshaping Arrays}
        You can manipulate the shape of an existing array using the \texttt{reshape()} method:
        \begin{lstlisting}
reshaped_array = np.arange(6).reshape((2, 3)) # Output: [[0, 1, 2], [3, 4, 5]]
        \end{lstlisting}
        Key Point: The total number of elements must remain the same when reshaping.
    \end{block}
\end{frame}

\begin{frame}[fragile]
    \frametitle{Creating Arrays - Example and Summary}
    \begin{block}{Example: Create and Reshape}
        Here’s a complete example demonstrating array creation and reshaping:
        \begin{lstlisting}
import numpy as np
        
# Creating a 1D array of the first 6 integers
array_1d = np.arange(6)  # Output: [0, 1, 2, 3, 4, 5]
        
# Reshaping the array to a 2x3 structure
array_reshaped = array_1d.reshape((2, 3))  # Output: [[0, 1, 2], [3, 4, 5]]
        \end{lstlisting}
    \end{block}
    
    \begin{block}{Summary}
        \begin{itemize}
            \item NumPy arrays are essential for numerical data manipulation.
            \item Use \texttt{np.array()} for creating arrays from lists, and various functions to create specific types of arrays.
            \item Reshape arrays effectively for computational needs.
        \end{itemize}
    \end{block}
\end{frame}

\begin{frame}[fragile]
    \frametitle{Creating Arrays - Conclusion}
    Mastering the creation and manipulation of NumPy arrays is a fundamental skill for any data science or machine learning application. 
    Practice creating arrays using different methods to become proficient with this powerful library!
\end{frame}

\begin{frame}
    \frametitle{Array Operations}
    \begin{block}{Overview}
        NumPy provides extensive capabilities for performing operations on arrays, significantly enhancing data manipulation skills.
    \end{block}
\end{frame}

\begin{frame}[fragile]
    \frametitle{Arithmetic Operations}
    \begin{block}{Element-wise Operations}
        Arithmetic operations are performed element-wise. Compatible dimensions are required for arrays.
    \end{block}
    
    \begin{itemize}
        \item Addition \texttt{(+)}
        \item Subtraction \texttt{(-)}
        \item Multiplication \texttt{(*)}
        \item Division \texttt{(/)}
        \item Exponentiation \texttt{(**)}
    \end{itemize}
    
    \begin{block}{Example}
        \begin{lstlisting}[language=Python]
import numpy as np

# Create two arrays
a = np.array([1, 2, 3])
b = np.array([4, 5, 6])

# Perform operations
addition = a + b      # Output: array([5, 7, 9])
subtraction = a - b   # Output: array([-3, -3, -3])
multiplication = a * b # Output: array([4, 10, 18])
division = a / b      # Output: array([0.25, 0.4, 0.5])
exponentiation = a ** 2 # Output: array([1, 4, 9])
        \end{lstlisting}
    \end{block}
\end{frame}

\begin{frame}[fragile]
    \frametitle{Logical Operations}
    \begin{block}{Element-wise Comparisons}
        Logical operations return Boolean arrays useful for filtering data.
    \end{block}

    \begin{itemize}
        \item Equal \texttt{(==)}
        \item Not Equal \texttt{(!=)}
        \item Greater Than \texttt{(>)}
        \item Less Than \texttt{(<)}
        \item Logical And \texttt{(&)}
        \item Logical Or \texttt{(|)}
    \end{itemize}
    
    \begin{block}{Example}
        \begin{lstlisting}[language=Python]
# Create an array
a = np.array([1, 2, 3, 4, 5])

# Logical operations
greater_than_two = a > 2   # Output: array([False, False,  True,  True,  True])
even_numbers = a % 2 == 0   # Output: array([False,  True, False,  True, False])
combined_condition = (a > 2) & (a % 2 == 0) # Output: array([False, False, False,  True, False])
        \end{lstlisting}
    \end{block}
\end{frame}

\begin{frame}[fragile]
    \frametitle{Aggregation Functions}
    \begin{block}{Summarizing Data}
        NumPy provides functions for aggregating array data.
    \end{block}
    
    \begin{itemize}
        \item \texttt{np.sum()}: sum of all elements
        \item \texttt{np.mean()}: average of elements
        \item \texttt{np.max()} / \texttt{np.min()}: largest/smallest element
    \end{itemize}
    
    \begin{block}{Example}
        \begin{lstlisting}[language=Python]
a = np.array([1, 2, 3, 4, 5])

# Aggregation
total = np.sum(a)        # Output: 15
mean_value = np.mean(a)  # Output: 3.0
max_value = np.max(a)    # Output: 5
min_value = np.min(a)    # Output: 1
        \end{lstlisting}
    \end{block}
\end{frame}

\begin{frame}
    \frametitle{Summary and Next Steps}
    \begin{itemize}
        \item Arithmetic Operations enable element-wise calculations.
        \item Logical Operations allow for comparisons producing Boolean results.
        \item Aggregation Functions help summarize and analyze data.
    \end{itemize}
    \begin{block}{Next Steps}
        We will explore \textbf{Indexing and Slicing} in NumPy, allowing us to access and modify specific elements within arrays.
    \end{block}
\end{frame}

\begin{frame}[fragile]
    \frametitle{Indexing and Slicing - Introduction}
    \begin{block}{Overview}
        In this section, we will learn how to access and modify specific elements within NumPy arrays using indexing and slicing techniques.
    \end{block}
\end{frame}

\begin{frame}[fragile]
    \frametitle{Indexing in NumPy}
    \begin{itemize}
        \item \textbf{Understanding Indexing:} Access elements using zero-based indexing.
        \item Key Points:
            \begin{itemize}
                \item \textbf{1D Arrays:} Use a single index.
                \item \textbf{2D Arrays:} Use a tuple of indices.
            \end{itemize}
    \end{itemize}
    \begin{block}{Example}
        \begin{lstlisting}[language=Python]
import numpy as np
arr = np.array([1, 2, 3, 4, 5])
print(arr[0])  # Outputs: 1

matrix = np.array([[1, 2, 3], [4, 5, 6]])
print(matrix[1, 2])  # Outputs: 6
        \end{lstlisting}
    \end{block}
\end{frame}

\begin{frame}[fragile]
    \frametitle{Modifying Elements and Slicing}
    \begin{block}{Modifying Elements}
        Modify elements using indexing:
        \begin{lstlisting}[language=Python]
arr[2] = 10
print(arr)  # Outputs: [1 2 10 4 5]
        \end{lstlisting}
    \end{block}
    \begin{block}{Understanding Slicing}
        Access a range of elements using the syntax: \texttt{array[start:end]}
        \begin{itemize}
            \item \textbf{Excludes:} The end index.
            \item \textbf{1D Array Example:}
        \end{itemize}
        \begin{lstlisting}[language=Python]
print(arr[1:4])  # Outputs: [2 10 4]
        \end{lstlisting}
    \end{block}
\end{frame}

\begin{frame}[fragile]
    \frametitle{Advanced Slicing Techniques}
    \begin{itemize}
        \item \textbf{Step (Stride):} Specify a step size.
        \begin{lstlisting}[language=Python]
print(arr[::2])  # Outputs: [1 30 5]
        \end{lstlisting}
        \item \textbf{Negative Indices:} Access elements from the end.
        \begin{lstlisting}[language=Python]
print(arr[-1])  # Outputs: 5
        \end{lstlisting}
    \end{itemize}
    \begin{block}{Summary}
        \begin{itemize}
            \item Indexing for single elements, slicing for a range.
            \item Key for effective data manipulation in NumPy.
        \end{itemize}
    \end{block}
\end{frame}

\begin{frame}[fragile]
    \frametitle{Practice Example}
    \begin{block}{Task}
        Create a NumPy array with values from 1 to 10. Access and modify elements using both indexing and slicing. What results do you get?
    \end{block}
    \begin{block}{Conclusion}
        Hands-on practice will solidify your understanding of indexing and slicing in NumPy!
    \end{block}
\end{frame}

\begin{frame}
    \frametitle{Statistical Functions}
    \begin{block}{Overview}
        NumPy is a powerful library in Python primarily used for numerical computing. 
        One of its key features is the availability of built-in statistical functions that simplify data analysis, enabling users to perform statistical operations on large datasets.
    \end{block}
\end{frame}

\begin{frame}[fragile]
    \frametitle{Key Statistical Functions - Part 1}
    \begin{enumerate}
        \item \textbf{Mean}
        \begin{itemize}
            \item \textbf{Function}: \texttt{numpy.mean()}
            \item \textbf{Description}: Computes the average of array elements.
            \item \textbf{Example}:
            \begin{lstlisting}
import numpy as np
data = [1, 2, 3, 4, 5]
mean_value = np.mean(data)  # Output: 3.0
            \end{lstlisting}
        \end{itemize}
        
        \item \textbf{Median}
        \begin{itemize}
            \item \textbf{Function}: \texttt{numpy.median()}
            \item \textbf{Description}: Finds the middle value in an array.
            \item \textbf{Example}:
            \begin{lstlisting}
median_value = np.median(data)  # Output: 3.0
            \end{lstlisting}
        \end{itemize}
    \end{enumerate}
\end{frame}

\begin{frame}[fragile]
    \frametitle{Key Statistical Functions - Part 2}
    \begin{enumerate}
        \setcounter{enumi}{2} % Continue the enumeration
        \item \textbf{Variance}
        \begin{itemize}
            \item \textbf{Function}: \texttt{numpy.var()}
            \item \textbf{Description}: Measures the spread of data points around the mean.
            \item \textbf{Example}:
            \begin{lstlisting}
variance_value = np.var(data)  # Output: 2.0
            \end{lstlisting}
        \end{itemize}
        
        \item \textbf{Standard Deviation}
        \begin{itemize}
            \item \textbf{Function}: \texttt{numpy.std()}
            \item \textbf{Description}: Provides a measure of variation of set values.
            \item \textbf{Example}:
            \begin{lstlisting}
std_dev_value = np.std(data)  # Output: 1.4142135623730951
            \end{lstlisting}
        \end{itemize}
        
        \item \textbf{Correlation Coefficient}
        \begin{itemize}
            \item \textbf{Function}: \texttt{numpy.corrcoef()}
            \item \textbf{Description}: Returns the correlation matrix of the input data.
            \item \textbf{Example}:
            \begin{lstlisting}
x = np.array([1, 2, 3])
y = np.array([4, 5, 6])
correlation = np.corrcoef(x, y)  # Output: [[1. 1.], [1. 1.]]
            \end{lstlisting}
        \end{itemize}
    \end{enumerate}
\end{frame}

\begin{frame}[fragile]
    \frametitle{Key Points and Conclusion}
    \begin{block}{Key Points to Emphasize}
        \begin{itemize}
            \item \textbf{Ease of Use}: Rapid statistical calculations on multi-dimensional arrays.
            \item \textbf{Performance}: Optimized implementations for large datasets.
            \item \textbf{Integration}: Works seamlessly with other libraries for comprehensive data analysis workflows.
        \end{itemize}
    \end{block}
    
    \begin{block}{Conclusion}
        Built-in statistical functions in NumPy are essential tools for data analysis, providing straightforward solutions to complex statistical calculations. Familiarity with these functions enhances your capability to derive insights from data efficiently.
    \end{block}
\end{frame}

\begin{frame}[fragile]
    \frametitle{Integrating NumPy with Other Libraries - Overview}
    \begin{block}{Brief Summary}
        NumPy is the foundational library for numerical computing in Python. Its integration with pandas and Matplotlib enhances data manipulation and visualization.
    \end{block}
\end{frame}

\begin{frame}[fragile]
    \frametitle{Integrating NumPy with Pandas}
    \begin{itemize}
        \item \textbf{Pandas}:
        \begin{itemize}
            \item A data analysis library utilizing NumPy arrays.
            \item Provides structures like DataFrames and Series.
        \end{itemize}

        \item \textbf{Integration}:
        \begin{itemize}
            \item DataFrames are built on top of NumPy arrays, allowing use of fast array operations.
        \end{itemize}
    \end{itemize}

    \begin{block}{Example Code}
    \begin{lstlisting}[language=Python]
import pandas as pd
import numpy as np

# Creating a NumPy array
data = np.array([[1, 2, 3], [4, 5, 6]])

# Creating a DataFrame from a NumPy array
df = pd.DataFrame(data, columns=['A', 'B', 'C'])
print(df)
    \end{lstlisting}
    \end{block}
\end{frame}

\begin{frame}[fragile]
    \frametitle{Integrating NumPy with Matplotlib}
    \begin{itemize}
        \item \textbf{Matplotlib}:
        \begin{itemize}
            \item A library for static, animated, and interactive visualizations.
        \end{itemize}

        \item \textbf{Integration}:
        \begin{itemize}
            \item Allows direct plotting from NumPy arrays.
        \end{itemize}
    \end{itemize}

    \begin{block}{Example Code}
    \begin{lstlisting}[language=Python]
import matplotlib.pyplot as plt

# Creating a NumPy array for x and y values
x = np.linspace(0, 10, 100)
y = np.sin(x)

# Plotting using Matplotlib
plt.plot(x, y)
plt.title("Sine Wave")
plt.xlabel("X-axis")
plt.ylabel("Y-axis")
plt.grid()
plt.show()
    \end{lstlisting}
    \end{block}
\end{frame}

\begin{frame}
    \frametitle{Hands-On Lab Session}
    \begin{block}{Overview of Exercises Using Real Data Sets}
        In this lab session, we will explore practical applications of NumPy through hands-on exercises utilizing real-world data sets. 
        The focus will be on implementing key NumPy functions and features to manipulate and analyze numeric data effectively.
    \end{block}
\end{frame}

\begin{frame}
    \frametitle{Objectives}
    \begin{itemize}
        \item Understand how to work with NumPy arrays in real-world contexts.
        \item Perform data manipulation tasks including selection, filtering, and basic statistical analysis.
        \item Integrate NumPy with real data sets for practical applications.
    \end{itemize}
\end{frame}

\begin{frame}[fragile]
    \frametitle{Exercise Breakdown}
    \begin{enumerate}
        \item \textbf{Data Set Loading and Inspection}
        \begin{itemize}
            \item Load a sample data set (e.g., global temperatures) in CSV format.
            \begin{lstlisting}
import numpy as np
import pandas as pd

# Load the dataset
data = pd.read_csv('data/global_temperatures.csv')
# Convert to NumPy array
temperatures = data['AvgTemperature'].to_numpy()
print(temperatures)
            \end{lstlisting}
            \item \textbf{Key Point:} Familiarize yourself with data types and structures.
        \end{itemize}
        
        \item \textbf{Basic Array Operations}
        \begin{itemize}
            \item Perform operations such as reshaping, slicing, and indexing.
            \item Extract the last 10 entries of temperature data.
            \begin{lstlisting}
last_ten = temperatures[-10:]
print(last_ten)
            \end{lstlisting}
            \item \textbf{Key Point:} Understand indexing to focus on specific parts of your data.
        \end{itemize}
    \end{enumerate}
\end{frame}

\begin{frame}[fragile]
    \frametitle{Exercise Breakdown (Contd.)}
    \begin{enumerate}[resume]
        \item \textbf{Statistical Analysis}
        \begin{itemize}
            \item Compute statistical measures: mean, median, standard deviation.
            \begin{lstlisting}
mean_temp = np.mean(temperatures)
median_temp = np.median(temperatures)
std_temp = np.std(temperatures)
print(f'Mean: {mean_temp}, Median: {median_temp}, Std Dev: {std_temp}')
            \end{lstlisting}
            \item \textbf{Illustration:} Graphically interpret the data distribution.
        \end{itemize}
        
        \item \textbf{Filtering and Conditional Operations}
        \begin{itemize}
            \item Filter data based on conditions (e.g., temperatures above a certain threshold).
            \begin{lstlisting}
high_temps = temperatures[temperatures > 30]
print(high_temps)
            \end{lstlisting}
            \item \textbf{Key Point:} Use conditional filtering to analyze trends or outliers.
        \end{itemize}
    \end{enumerate}
\end{frame}

\begin{frame}[fragile]
    \frametitle{Integrating with Other Libraries}
    \begin{enumerate}
        \item \textbf{Visualization}
        \begin{itemize}
            \item Use Matplotlib to visualize data extracted through NumPy.
            \begin{lstlisting}
import matplotlib.pyplot as plt

plt.plot(temperatures)
plt.title('Average Temperatures Over Time')
plt.xlabel('Time')
plt.ylabel('Temperature')
plt.show()
            \end{lstlisting}
        \end{itemize}
        
        \item \textbf{Questions for Reflection}
        \begin{itemize}
            \item How do NumPy's array operations simplify data manipulation compared to standard Python lists?
            \item What insights can you derive from the statistical analysis and visualizations of the data?
        \end{itemize}
    \end{enumerate}
\end{frame}

\begin{frame}
    \frametitle{Next Steps}
    As we progress to the next slide, we will summarize our key learnings and share additional resources for continued exploration of NumPy and its capabilities.
    
    This lab session is designed to immerse you in hands-on applications of NumPy, reinforcing your understanding of how to manipulate and analyze data effectively. Prepare to engage with the data and ask questions that deepen your comprehension of the material!
\end{frame}

\begin{frame}[fragile]
    \frametitle{Conclusion and Resources - Key Takeaways}
    
    \begin{enumerate}
        \item \textbf{What is NumPy?} \\
        NumPy (Numerical Python) is an open-source library for numerical computing in Python that supports arrays, matrices, and a rich set of mathematical functions.
        
        \item \textbf{Core Features:}
        \begin{itemize}
            \item \textbf{N-Dimensional Arrays:} Efficient manipulation of large datasets.
            \item \textbf{Broadcasting:} Facilitates arithmetic operations on differently shaped arrays.
            \item \textbf{Vectorization:} Element-wise operations without explicit loops for better performance.
        \end{itemize}
    \end{enumerate}
\end{frame}

\begin{frame}[fragile]
    \frametitle{Conclusion and Resources - Basic Operations}
    
    \begin{enumerate}
        \item \textbf{Creating Arrays:}
        \begin{lstlisting}[language=Python]
import numpy as np
array_1d = np.array([1, 2, 3])
array_2d = np.array([[1, 2, 3], [4, 5, 6]])
        \end{lstlisting}
        
        \item \textbf{Array Operations:}
        \begin{itemize}
            \item \textbf{Addition:}
            \begin{lstlisting}[language=Python]
result = array_1d + 10  # Adds 10 to each element
            \end{lstlisting}
            \item \textbf{Dot Product:}
            \begin{lstlisting}[language=Python]
dot_result = np.dot(array_2d, array_1d)  # Matrix multiplication
            \end{lstlisting}
        \end{itemize}
    \end{enumerate}
\end{frame}

\begin{frame}[fragile]
    \frametitle{Conclusion and Resources - Additional Learning Resources}
    
    \begin{block}{Useful Resources}
        \begin{itemize}
            \item \textbf{Official Documentation:} \\
            \texttt{https://numpy.org/doc/stable/}
            
            \item \textbf{Online Tutorials:} \\
            \texttt{https://www.w3schools.com/python/numpy_intro.asp}
            
            \item \textbf{Books:}
            \begin{itemize}
                \item "Python for Data Analysis" by Wes McKinney
                \item "Numerical Methods in Engineering with Python"
            \end{itemize}
            
            \item \textbf{Online Courses:}
            \begin{itemize}
                \item Coursera: Courses on Data Science & Machine Learning
                \item edX: Programming courses using NumPy
            \end{itemize}
        \end{itemize}
    \end{block}
\end{frame}


\end{document}